\chapter{细胞生物学绪论}

\section{细胞生物学研究的内容与现状}

\subsection{细胞生物学是显带生命科学的重要基础学科}

\firstgls{cellbiology}细胞生物学是研究细胞基本生命活动规律的科学，它主要从不同层次（显微、亚显微与分子水平）上研究细胞结构与功能，细胞增殖、分化、衰老与死亡，细胞信号转导，细胞基因的表达与调控，细胞的起源与进化等。分子细胞生物学是当今细胞生物学研究的重点，细胞工程及与之相关的组织工程和修复医学是21世纪生物工程发展的重要组成部分。可以看到，细胞的结构与基本生命活动的研究已经越来越深入，并已经成为21世纪初生命科学研究的重要领域之一。

生命是多层次、非线性、多侧面的复杂结构体系，而细胞是生命体结构与生命活动的基本单位，有了细胞才有完整的生命活动。细胞的研究是生命科学的基础，也是现代生命科学发展的重要支柱。早在1925年，生物学大师Edmund Beecher Wilson就提出：“一切生命的关键问题都要到细胞中去寻找”（The key to every biological problem must finally be sought in the cell）。重温这句名言，至今仍感其内涵深刻。

生物的生殖发育、遗传、神经（脑）活动等重大生命现象的研究都要以细胞为基础。多细胞生物的生长发育是依靠细胞增殖、\firstgls{cell_differentiation} 与\firstgls{cell_apoptosis}来实现的。人脑的活动是靠$ 10^{{12}} $个细胞及其相互协调而进行的。一切疾病发病机制也要以细胞病变的研究为基础。以基因工程和蛋白质工程为核心的显带生物技术主要是通过以细胞操作为基础而进行的。因此，细胞生物学与农业、医学、生物技术的发展有密不可分的关系，它将在解决人类面临的重大问题、促进经济和社会发展中发挥重要的基础作用。

在我国基础科学发展规划中，细胞生物学、分子生物学、神经生物学和生态学被并列为生命科学的四大基础学科，反映了显带生命科学的发展趋势。

由于分子生物学概念、方法和技术的引入，细胞生物学在近些年取得了突破性的进展，产生了许多新的生长点，并逐渐形成新的概念与新的领域。很多科学家认为，在21世纪，细胞生物学将继续徐梦发展，因为它是解释生命奥秘不可缺少的``主角''。细胞生物学自身优势多层次的复杂结构体系，它是物质（结构）、信息与能量相互``辉映''的综合体。它的很多基本生命活动过程是如何有序而被自动调控的还不是很清楚，生物大分子如何逐级有序地组装成行使生命活动的细胞基本结构体系页相当朦胧。因此，细胞的基础研究势必影响21世纪生命科学的整体发展速度。


可以概况地说，细胞生物学是应用显带物理学与化学的技术成就和分子生物学的概念与方法，以细胞作为生命活动的基本单位为出发点，探索生命活动规律的学科，其核心问题是将遗传与发育在细胞水平上结合起来。还应指出，细胞生物学是一门迅速发展中的新兴学科，其研究内容与范畴往往与生命科学的其他学科交错在一起，甚至目前很难为细胞生物学划出一个明确的范围。

\subsection{细胞生物学的主要研究内容}

细胞生物学研究与教学内容一般可分为细胞结构与功能和细胞重要生命活动两大基本部分，但它们又很难割裂开。从20世纪60年代开始，细胞超微结构研究积累的大量资料，大大充实与拓宽了细胞结构与功能的知识范畴，在细胞生物学教科书中，与细胞结构和功能相关的知识所占比例较多。从70年代中期开始，由于分子生物学概念、内容与方法的引入，细胞生物学面貌发生了深刻的变化，不仅使细胞结构和功能的研究更深入，对细胞重大生命活动规律及其调控机制的研究也取得了巨大进展，极大地丰富与改变了细胞生物学的知识结构。因此，现代细胞生物学教科书中细胞重要生命活动的知识所占比例也越来越大。

当前细胞生物学研究内容大致归纳为以下领域：

\subsubsection{细胞核、染色体以及基因表达的研究}

细胞核是遗传物质DNA储存的场所，也是遗传信息进行转录的场所。染色质与染色体是遗传物质的载体，核仁是转录rRNA以及组装核糖体亚单位的具体场所，核孔复合体是物质之间交换与信息交流的结构。细胞核与染色体的研究历来是经典细胞学的重点，也是细胞遗传学的核心部分。而现代细胞生物学的核心课题之一就是研究染色体结构动态变化与基因表达及其调控的关系，它是目前细胞生物学、遗传学与发育生物学在细胞水平与分子水平上相结合的最活跃的热门课题之一，也是后基因组时代生命科学主要的研究内容之一。

\subsubsection{生物膜与细胞器的研究}

生物膜是细胞结构的重要基础，细胞质膜出现之后，才有了真正意义上的生命体。大部分细胞器（包括细胞核）也是以生物膜为基础构建的。近年来，生物膜研究的主要内容是\underline{膜的结构模型}、\underline{物质的跨膜运输}及\underline{信息跨膜传递的机制}。磷脂双分子层与膜蛋白的相互关系是研究生物膜结构与功能的重要内容。几年来，在膜的识别与受体效应、蛋白质分子跨膜运输与定向分选等方面取得了巨大进展。

细胞器的研究历来是认识细胞结构与功能的重要组成部分。线粒体与叶绿体结构与换能机制的研究已很深入，线粒体DNA与叶绿体DNA的发现及其半自主性的研究是人们对这两种细胞器又有了新的认识。近年来对内质网、高尔基体与溶酶体功能的研究也增添了许多新的知识。核糖体RNA能催化肽链的合成的发现，大大加深了人们对核糖体功能的认识，同时也显示了核糖体在生命起源与进化中的重要地位。

\subsubsection{细胞骨架体系的研究}

细胞骨架体系的研究在细胞生物学中是一个重要的研究领域。广义的细胞骨架概念应该包括\underline{细胞质骨架}与\underline{核骨架}两大部分。

细胞骨架体系的研究越来越受到重视，因为细胞骨架在维持细胞形态与保持细胞内部结构的合理布局中起主要作用。细胞骨架与一系列重要生命活动，诸如细胞内大分子的运输与细胞器的运动、细胞信息的传递、基因表达与大分子加工等均有密切关系。今年来，细胞核骨架的研究是颇引人关注的领域之一。核骨架包括\underline{核基质}、\underline{核纤层}和\underline{核孔复合体}，不仅参与核染色体的构建，而且与基因表达关系密切。

\subsubsection{细胞增殖及其调控}

一切动植物的生长及发育都是通过细胞的增殖与分化来实现的。研究细胞增殖的基本规律及其调控机制，不仅是控制生物生长及发育的基础，而且是研究癌变发生及逆转的重要途径。目前国际上研究细胞增殖的调控主要从两方面进行：

\begin{enumerate}
	\item 从环境或有机体中寻找控制细胞增殖的因子，以及阐明它们的作用机制；各种生长因子的发现及其作用机制的揭示是这一领域中重要的进展。
	\item 寻找控制细胞增殖的关键性基因，并通过调节基因产物来控制细胞的增殖。
\end{enumerate}

细胞的癌基因与抑癌基因及其表达产物均与细胞增殖有关。

\subsubsection{细胞分化及其调控}

细胞分化是生物发育的基础。近年，细胞分化的研究已越来越显示出其重要性，也是细胞生物学、发育生物学与遗传学的重要汇合点。近年生物科学的发展，尤其是分子生物学技术的建立，已为细胞分化机制的研究奠定了良好的基础，这也是近年发育生物学蓬勃发展的重要原因。

一个受精卵通过分裂与分化如何发育为复杂的有机体，是生命科学中引人入胜的课题之一。利用哺乳动物的体细胞进行动物克隆，进一步揭示了细胞的``全能性''，使人们认识到可以控制细胞的分化，而且可以将已分化的细胞进行``去分化''，并使其分裂与再分化，克隆出新的个体，这为控制生物的生长发育展示了诱人的前景。

目前认为\underline{细胞分化的本质是细胞内基因选择性表达特异功能蛋白质的过程}。近年，细胞分化及其调控机制的研究主要集中在\underline{编码特异蛋白质的基因的选择性表达规律及其调控}方面。围绕干细胞的研究及其再组织工程中的应用前景，极大地促进了细胞分化领域的研究。

\subsubsection{细胞的衰老与凋亡}

细胞衰老的研究是研究人与动植物寿命的基础，细胞总体的衰老导致个体的衰老，但细胞的衰老与有机物的衰老又是不同的概念。不少科学家认为细胞衰老的研究将成为21世纪的热门课题之一。

目前多数科学家是用细胞体外培养的方法来研究细胞衰老的规律。大量实验说明，动物二倍体细胞在体外分裂与传代的次数是有限的，从而推测体内细胞的寿命受分裂次数的限制，细胞的衰老是必然的规律。人们总是希望通过细胞衰老因素与因子的研究来延长细胞的寿命。近年人们力图寻找细胞中的衰老基因及相关信号转导途径。

细胞凋亡的研究是生命科学中发展起来的重要的新兴领域之一。细胞凋亡是由一系列基因控制并受复杂的信号调节的细胞自然死亡现象，细胞凋亡可能是生物正常生理发育与病理过程中的重要平衡因素。

\subsubsection{细胞的起源与进化}

细胞的起源与进化的研究是重要的理论问题，也是难度很大的研究课题，我们应该十分尊重先驱科学家在这一领域所取得的成果。

根据古微生物的证据，原始细胞大约在35亿年，或更早到38亿年前就在地球上出现。生命物质只有当其以细胞形式出现时才能稳定地在地球上存在下去，但是这个进程已成为历史，在现在实验条件下模拟30多亿年前的地球环境，研究细胞的起源过程几乎是不大可能的。目前有关细胞起源学说在很大程度上是推理性的，还需要补充更多的事实证据。

关于真核细胞的起源与进化已取得了较多的成果，近年由于分子生物学方法引入生物进化的研究，为细胞进化与细胞器进化的研究增添了很多新的内容。

\subsubsection{细胞工程}

细胞工程是细胞生物学与遗传学的交叉领域，这种改造细胞的技术是生物工程技术的重要组成部分。它不仅对工农业生成和医学实践有重要意义，而且也是认识细胞生命活动规律的一种重要途径与手段。

细胞工程能用人工方法使不同种细胞的基因或基因组重组到杂交细胞中，或者使基因与基因组由一种细胞转移到另一种细胞中，并使越过种的障碍的基因成为可能。由此人们开始探索人工创造新的遗传性状的细胞。动植物体细胞的杂交试验一直使细胞工程中最活跃的领域。通过动物体细胞杂交而获得单克隆抗体技术的建立，是细胞工程中最富成果性的工作范例。近年，在世界范围内兴起的用哺乳动物体细胞克隆而获得无性繁殖胚胎与个体，是细胞工程最引人注意的进展之一。人胚胎干细胞系的建立，为细胞工程中的一个新的研究领域---干细胞工程的研究与应用打下基础。

还应该说明，当前细胞生物学研究的范畴远不止上述的内容，其他如细胞信号转导已成为细胞生物学最热门的领域之一，一切重大生命活动都涉及到细胞信号转导与调控。此外，细胞外基质、细胞社会学与细胞免疫学等研究，近年也由较快的发展。

\subsection{当前细胞生物学研究的总趋势与重点领域}

细胞生物学与分子生物学相互渗透与交融是总的发展趋势。无论是细胞结构与功能的深入研究，还是对细胞重大生命活动规律的探索，都需要用分子生物学的新概念与新方法，在分子水平上进行研究。换句话说，细胞分子生物学或分子细胞生物学将是今后相当一段时间的主流学科方向。

\subsubsection{当前细胞生物学研究中的几大基本问题}

当前细胞生物学研究的课题归纳起来包括以下几个根本性的问题：

\paragraph{基因组在细胞内是如何在时间与空间上有序表达的？}
如果将一些基因放在试管内，只要条件都满足，它们都可以完成表达。但是在细胞内环境中，其表达程度都将受到严格的调节与控制，它们能否表达？这就是细胞作为结构与生命活动的基本单位的奥妙所在。细胞结构与生命活动的有序性是十分复杂的，是非线性调控过程。这种调控网络可能是迄今任何一台计算机都无法比拟的。

\paragraph{基因表达的产物，主要是结构蛋白与核酸、脂质、多糖及其复合物，如何逐级组装成能行使生命活动的基本结构体系及各种细胞器？}
这种自组装过程的调控程序与调控机制是什么？这又是当前极富挑战性的领域。我们认为它应是新兴的结构生物学的重要组成部分，生命活动很多本质问题都能在这里找到答案。

\paragraph{基因表达的产物，主要是大量活性因子与信号分子，是如何调节诸如细胞的增殖、分化、衰老、与凋亡等细胞最重要的生命活动过程的？}
这些研究工作现在方兴未艾。

\subsubsection{当前细胞基本生命活动研究的若干重大课题}

\paragraph{染色体DNA与蛋白质相互作用关系---主要是非组蛋白对基因组的作用。}
近年来染色体（质）的概念有很大的改变。可以认为染色体是一条巨大的DNA分子与其结合蛋白形成的动态复合结构。其存在状态可以分为细胞间期染色质与分裂期染色体。两者在细胞周期中交替与互相转换，主要是由DNA与蛋白质动态结合关系所决定的。DNA是遗传密码的分子载体，蛋白质是调节DNA复制与转录活性并参与染色体构建的元件。染色体DNA结合蛋白可分为组蛋白与非组蛋白。核小体的发现及其结构的阐明，对5种组蛋白的作用已基本清楚。大量与种类繁多的非组蛋白与DNA的相互关系是研究的主要问题，特别是它们对DNA的复制、基因组有序表达与染色体高层次构建及动态结构变化的作用是核心问题。同时，DNA分子与组蛋白的修饰在基因表达调控中的作用，即表观遗传学的研究也越来越受到人们的关注。

\paragraph{细胞增殖、分化、凋亡（程序性死亡）的相互关系及其调控}
细胞增殖、分化、凋亡与衰老是细胞最重要的生命活动现象，它们既相互联系又相互制约，并收到一系列基因的调控。一切动植物生长发育都是依靠细胞增殖、分化与凋亡来实现的。细胞癌变往往是由于细胞增殖、分化或凋亡失控所致。受精卵如何有条不紊地发育为复杂的成体，一直是人们企图揭开的奥秘。由于近年来遗传与发育在分子水平上的统一并取得了巨大进展，为揭示这一问题提供了光明前景。而细胞分化及其调控机制的研究，将是解决这一重大问题的基础。

细胞凋亡的研究是近年生命科学中发展起来的热点课题。细胞凋亡是受外环境多种信号刺激，由基因调控的``主动自然''死亡，在细胞增殖、分化及多种病理过程中具有重要作用。细胞凋亡可能是发育进程的重要平衡因素。例如，哺乳动物脑的发育过程中有$20 \%\sim 50\% $的脑神经元发生凋亡；泌乳细胞在婴儿断乳后很快凋亡，代之以脂肪细胞。应该凋亡的细胞不凋亡也是致癌的重要因素。

人们已经通过促进细胞凋亡来为癌症治疗提供依据及手段。植物细胞与单细胞生物有无类似动物细胞的凋亡现象也是我们要探索的问题。不论是细胞增殖与分化还是细胞凋亡的研究，核心问题依然是其调控机制的研究。

\paragraph{细胞信号转导的研究}
细胞信号转导的研究与生命科学中的许多重要问题密切相关，它已成为了了解错综复杂的生命现象的主要切入点，并成为生命科学多领域、多层次联系的纽带。细胞信号转导的研究对了解生命本质与细胞基本生命活动有重要理论意义。

当前细胞信号转导的基本内容是：\textbf{细胞间信号传递}：重点是信号分子的结构与功能、信号分子与受体相互作用机制；\textbf{受体与信号跨膜转导}：G蛋白的发现与较深入的研究使受体转换器的概念与结构功能更具体化；\textbf{细胞内信号传递途径}：蛋白磷酸化与去磷酸化可能使信号转导的关键。把信号转导与基因表达调控联系起来的研究可能使最吸引人的课题之一。

应该指出，细胞信号转导使通过复杂的网络系统实现的，呈现高度的非线性、多侧面的特点。细胞信号转导的基础使蛋白质与蛋白质的相互作用，特别使功能蛋白质的相互做哟用。信号转导通路可能不是单一的或简单的信息叠加。因此，研究细胞信号转导通路质检的相互作用将成为该课题的热点。

\paragraph{细胞结构体系的组装}

生物大分子如何逐级组装并最终形成赖以进行生命活动的细胞结构体系？这是当前生命科学面临的最基本问题之一。这一课题不仅是揭示细胞生命本质与奥妙的重要组成部分，而且也是深入研究细胞结构与功能的新的切入点。大分子嗯那个逐级自组装成能行使生命功能的基本细胞体系，本身就是引人入胜的重大问题。

在分子水平上研究与解释生命现象的本质固然十分重要，但细胞结构体系的自组装与去组装动态过程不能简单归属于分子水平的活动，尚需在更高层次上进行探索与研究。生命活动的基础是细胞内高度有序的、动态的复合结构体系，在亚细胞水平我们可以将细胞结构大致归纳为几大基本结构体系：\underline{由蛋白质与核酸构建的遗传信息结构体系}（包括染色体、核仁、核糖体等）、\underline{由蛋白质与脂质构成的膜结构体系}（包括细胞膜、核膜以及其他各种细胞器膜）、\underline{由蛋白质与蛋白质构成的细胞骨架结构体系}。由这些基本结构体系形成的细胞器与各级结构在细胞内进行组装、去组装、重组装与重建（改建）等过程。有不少细胞结构可以在体外进行组装与重组装。体外组装实验模式的建立，为研究细胞结构体系组装机制提供了很好的基础与手段。近年来，在这一领域进展最快的是细胞核体外组装（重建）、细胞骨架组装、人工染色体与人工膜组装等的研究。\\


最后必须指出，上述诸领域与课题不能概括细胞基本生命活动的所有重大问题。还有很多方面，诸如蛋白质合成、分选与跨膜定向运输、真核细胞起源与进化等，也对细胞基本生命活动规律的认识很重要。

美国国立卫生研究院（NIH，National Institutes of Health）在1988年底发表的一份题为《什么是当前科研领域热门话题？》（``What is popular in research today?''）的调查报告中指出，目前全球研究最热门的几种疾病：癌症（Cancer）、心血管病（Cardiovascular Diseases）、艾滋病（AIDS）和肝炎（Hepatitis）等传染病（Infectious Diseases）。研究最热门的方向是：细胞周期调控（Cell Cycle Control）、细胞凋亡（Cell Apoptosis）、细胞衰老（Cell Senescence）、细胞信号转导（Signal Transduction）和DNA的损伤与修复（DNA Damage and Repair）。2004年，NIH提出的生物医学发展规划中，更是把细胞信号转导网络作为今后10年研究的主要内容。

由于细胞是生命活动的基本单位，又是生命的缩影，它不仅体现了生命的多样性和统一性，更体现了生命的复杂性。显然，我们对细胞重大生命活动的认识，必须建立在创新思维的基础上，着力探索它们彼此之间的关系，以及这些重大生命活动与求解人类所面临的诸多问题之间的关系。为此，我们给出一个有关细胞重大生命活动及其相互关系的图示，可能有助于思考和研究这些问题。

\begin{center}
\pgfdeclarelayer{backwrapping}
\pgfdeclarelayer{frontwrapping}
\pgfsetlayers{backwrapping,main,frontwrapping}

\tikzset{
    wrapping/.style={
        draw=black!60, 
        line cap=round, 
        line join=round, 
        ultra thick},
    nucleosome/.style={
        fill=red!40, 
        fill opacity=.9, 
        draw=none},
    top cylinder/.style={
        fill=red!60, 
        fill opacity=.9
    }
}

\begin{tikzpicture}

\foreach \q [remember=\q as \p] in {1, 2, 3, 4}{
    \begin{scope}[shift={(\q*3,0)}, rotate=30]
        \path [nucleosome]
            (0,1)
            arc (90:270:0.375 and 1) -- (1.25,-1)
            arc (270:90:0.375 and 1) -- cycle;
        \path [top cylinder] 
            (1.625, 0) arc (0:360:0.375 and 1) -- cycle;
          
          \begin{scope}[shift={(0.25,0)}]
            \begin{pgfonlayer}{backwrapping}
                \draw [wrapping] 
                    (0.25, -1.125) 
                    \foreach \i in {180,185,...,360}
                        { -- (\i/720+0.375*sin -\i, 1.125*cos \i)};
                    \draw [wrapping] 
                        (0, 1.125)
                        arc(90:0:0.125cm and 0.25cm)
                        arc(0:-90:1cm and 1cm) coordinate (wrapping-start-\q);
            \end{pgfonlayer}
            
            \begin{pgfonlayer}{frontwrapping}
                \draw [wrapping] 
                    (0, 1.125) 
                    \foreach \i in {0,5,...,180}
                        { -- (\i/720+0.375*sin -\i,1.125*cos \i)}
                        [shift={(0.5,0)}]
                        (0, 1.125) 
                    \foreach \i in {0,5,...,150}{ -- (\i/720+0.375*sin -\i,1.125*cos \i)} coordinate (wrapping-end-\q);
            \end{pgfonlayer}
          \end{scope}
          
      \ifnum\q>1
        \draw [wrapping] (wrapping-end-\p) .. controls ++(-60:0.5cm)  and ++(180:0.25cm) .. (wrapping-start-\q);
      \fi
      \ifnum\q=4
        \draw [wrapping] (wrapping-end-\q) arc (210:270:1cm and 0.75cm);
      \fi
      \end{scope}
      }
  \end{tikzpicture}
  
  \tikz \draw[thick,rounded corners=16pt]
  (0,0) -- (0,2) -- (1,3.25) -- (2,2) -- (2,0) -- (0,2) -- (2,2) -- (0,0) -- (2,0);

\end{center}

\section{细胞学与细胞生物学发展简史}

现在许多科学家认为，可以把生物科学的发展划分为3个阶段：19世纪以及更早的时期，是以形态描述为主的生物科学时期；20世纪的前半个世纪，主要是实验的生物学时期；20世纪五六十年代以来，以DNA双螺旋的发现与中心法则的建立，开始了进入精细定位与定量的生物学时期。细胞学的发展史和细胞生物学学科的建立与发展，大致符合这一历史进程。

现代细胞生物学已发展到恨到的水平，已非经典细胞学可以比拟，但遵循科学发展的规律，我们仍有必要简要重温一下细胞学与细胞生物学发展的历史，这将有助于我们加深对细胞生物学学科的形成及其历史条件的认识。

细胞学与细胞生物学发展的历史大致可以划分为以下几个阶段：

\subsection{细胞的发现}


\subsection{细胞学说的建立及其意义}

\subsection{细胞学的经典时期}

\subsection{实验细胞学与细胞学的分支及其发展}

\subsection{细胞生物学学科的形成与发展}

\subsection{细胞生物学的主要学术组织、学术刊物与教科书}









